\documentclass[12pt, a4paper]{article}
\usepackage[utf8]{inputenc}
\usepackage[T1]{fontenc}
\usepackage[turkish]{babel}
\usepackage{amsmath}
\usepackage{amsfonts}
\usepackage{amssymb}
\usepackage{geometry}
\geometry{margin=2.5cm}

\title{\textbf{Mühendislik Matematiği Ara Sınav Çözümleri}}
\author{Çözümler}
\date{}

\begin{document}

\maketitle

\section*{Soru 1: Çok Değişkenli En Küçük Kareler Yöntemi}

\textbf{Soru:} $\{x_i, y_i, z_i; i=1, n\}$ veri grubuna $z(x,y) = ax^2y^2 + bxy + cx + dy + e$ modelini çakıştırmak için gerekli denklem sistemini geliştiriniz.

\textbf{Çözüm:}

Amaç, model ile veri noktaları arasındaki hatanın kareleri toplamını ($E$) minimize etmektir.

Modelimiz: 
\[ z_{model} = ax^2y^2 + bxy + cx + dy + e \]

Hata Kareler Toplamı ($E$):
\begin{equation}
E = \sum_{i=1}^{n} \left[ z_i - (ax_i^2y_i^2 + bx_iy_i + cx_i + dy_i + e) \right]^2
\end{equation}

Bu ifadenin katsayılara ($a, b, c, d, e$) göre kısmi türevlerini alıp sıfıra eşitleriz:

\begin{align*}
\frac{\partial E}{\partial a} &= -2 \sum [z_i - (\dots)] (x_i^2y_i^2) = 0 \\
\frac{\partial E}{\partial b} &= -2 \sum [z_i - (\dots)] (x_iy_i) = 0 \\
\frac{\partial E}{\partial c} &= -2 \sum [z_i - (\dots)] (x_i) = 0 \\
\frac{\partial E}{\partial d} &= -2 \sum [z_i - (\dots)] (y_i) = 0 \\
\frac{\partial E}{\partial e} &= -2 \sum [z_i - (\dots)] (1) = 0
\end{align*}

Düzenlendiğinde elde edilen \textbf{Normal Denklemler Sistemi} (Matris formunda yazılmaya uygun hali):

\begin{equation}
\begin{aligned}
a\sum x^4y^4 + b\sum x^3y^3 + c\sum x^3y^2 + d\sum x^2y^3 + e\sum x^2y^2 &= \sum z x^2y^2 \\
a\sum x^3y^3 + b\sum x^2y^2 + c\sum x^2y + d\sum xy^2 + e\sum xy &= \sum z xy \\
a\sum x^3y^2 + b\sum x^2y + c\sum x^2 + d\sum xy + e\sum x &= \sum z x \\
a\sum x^2y^3 + b\sum xy^2 + c\sum xy + d\sum y^2 + e\sum y &= \sum z y \\
a\sum x^2y^2 + b\sum xy + c\sum x + d\sum y + e\cdot n &= \sum z
\end{aligned}
\end{equation}
\textit{(Not: İndisler sadelik için yazılmamıştır, tüm toplamlar $i=1$ den $n$'e kadardır.)}

\hrulefill

\section*{Soru 2: Dalga Denkleminin Analitik Çözümü}

\textbf{Soru:} $\frac{\partial^2 U}{\partial t^2} - c^2 \frac{\partial^2 U}{\partial x^2} = 0$ ($0 < x < h, t > 0$) denkleminin genel çözümünü analitik yolla bulunuz.

\textbf{Çözüm:}

Bu homojen dalga denklemini çözmek için \textbf{Değişkenlerine Ayırma Yöntemi} uygulanır.
Çözümün $U(x,t) = X(x)T(t)$ formunda olduğunu varsayalım.

\begin{equation}
X(x)T''(t) - c^2 X''(x)T(t) = 0 \Rightarrow \frac{T''(t)}{c^2 T(t)} = \frac{X''(x)}{X(x)} = k
\end{equation}

Salınım yapan (fiziksel) bir çözüm için ayırma sabiti negatif seçilmelidir: $k = -\lambda^2$.

Bu durumda iki adet Adi Diferansiyel Denklem elde edilir:
\begin{align}
X''(x) + \lambda^2 X(x) &= 0 \\
T''(t) + c^2 \lambda^2 T(t) &= 0
\end{align}

Bu denklemlerin genel çözümleri:
\[ X(x) = A \cos(\lambda x) + B \sin(\lambda x) \]
\[ T(t) = C \cos(c\lambda t) + D \sin(c\lambda t) \]

Denklemin \textbf{Genel Çözümü}:
\begin{equation}
U(x,t) = \left[ A \cos(\lambda x) + B \sin(\lambda x) \right] \left[ C \cos(c\lambda t) + D \sin(c\lambda t) \right]
\end{equation}

\textbf{Yorum:} Matematiksel olarak bu ifade trigonometrik serilerin çarpımıdır. Fiziksel olarak ise, uzayda ($X$) sinüzoidal bir şekle sahip olan ve zamanla ($T$) yine sinüzoidal olarak salınan \textbf{Duran Dalgaları (Standing Waves)} temsil eder.

\hrulefill

\section*{Soru 3: Doğrusal Olmayan Model Uydurma}

\textbf{Soru:} $y(x) = AB^x$ modelini En Küçük Kareler yöntemine göre çözünüz.

\textbf{Çözüm:}

Model lineer değildir. Önce logaritma alarak lineerize (doğrusal) hale getirmeliyiz.

\begin{equation}
\ln(y) = \ln(A B^x) \Rightarrow \underbrace{\ln(y)}_{Y} = \underbrace{\ln(A)}_{a_0} + \underbrace{\ln(B)}_{a_1} x
\end{equation}

Dönüşümler:
$Y_i = \ln(y_i)$, $a_0 = \ln(A)$, $a_1 = \ln(B)$.
Yeni model: $Y = a_0 + a_1 x$ (Bir doğru denklemi).

Hata kareler toplamı:
\[ E = \sum_{i=1}^{n} (Y_i - a_0 - a_1 x_i)^2 \]

Türevler alınıp sıfıra eşitlendiğinde ($ \frac{\partial E}{\partial a_0}=0, \frac{\partial E}{\partial a_1}=0 $) aşağıdaki matris sistemi elde edilir:

\begin{equation}
\begin{bmatrix}
n & \sum x_i \\
\sum x_i & \sum x_i^2
\end{bmatrix}
\begin{bmatrix}
a_0 \\
a_1
\end{bmatrix}
=
\begin{bmatrix}
\sum Y_i \\
\sum x_i Y_i
\end{bmatrix}
\end{equation}

Bu sistem çözülüp $a_0$ ve $a_1$ bulunduktan sonra asıl katsayılar:
\[ A = e^{a_0}, \quad B = e^{a_1} \]

\hrulefill

\section*{Soru 4: Diferansiyel Denklem Özel Çözümü}

\textbf{Soru:} $9y''(t) + 6y'(t) + y(t) = 0$; $y(0)=1, y'(0)=0$

\textbf{Çözüm:}

Karakteristik denklem:
\[ 9r^2 + 6r + 1 = 0 \Rightarrow (3r+1)^2 = 0 \]
Kökler: $r_1 = r_2 = -\frac{1}{3}$ (Çakışık kök).

Genel Çözüm:
\begin{equation}
y(t) = c_1 e^{-t/3} + c_2 t e^{-t/3}
\end{equation}

Sınır koşullarını uygulayalım:
1. $y(0) = 1$:
\[ c_1 e^0 + c_2(0) = 1 \Rightarrow \mathbf{c_1 = 1} \]

2. Türev alalım $y'(t)$:
\[ y'(t) = -\frac{1}{3}c_1 e^{-t/3} + c_2 \left( 1 \cdot e^{-t/3} - \frac{1}{3}t e^{-t/3} \right) \]
$c_1=1$ yerine yazıp $t=0$ verelim ($y'(0)=0$):
\[ 0 = -\frac{1}{3}(1) + c_2(1 - 0) \Rightarrow c_2 = \frac{1}{3} \]

\textbf{Özel Çözüm:}
\begin{equation}
y(t) = e^{-t/3} + \frac{1}{3} t e^{-t/3} = e^{-t/3} \left( 1 + \frac{t}{3} \right)
\end{equation}

\end{document}